\documentclass[12pt,a4paper]{article}

\usepackage[utf8]{inputenc}
\usepackage[T1]{fontenc}
\usepackage[margin=1in]{geometry}
\usepackage{amsmath,amssymb}
\usepackage{hyperref}
\usepackage{booktabs}
\usepackage{enumitem}
\usepackage{xcolor}
\usepackage{fancyhdr}
\usepackage{setspace}
\usepackage{colortbl}
\usepackage{multirow}
\usepackage{longtable}

\onehalfspacing
\pagestyle{fancy}
\fancyhf{}
\fancyhead[R]{\textit{\small Work Plan --- Spectral Binary Neural Networks}}
\fancyfoot[C]{\thepage}
\renewcommand{\headrulewidth}{0.4pt}

\definecolor{phase1}{HTML}{E8F5E9}
\definecolor{phase2}{HTML}{E3F2FD}
\definecolor{phase3}{HTML}{FFF3E0}
\definecolor{phase4}{HTML}{F3E5F5}
\definecolor{phase5}{HTML}{FFEBEE}
\definecolor{reviewfix}{HTML}{FFF9C4}

\newcommand{\review}{\texorpdfstring{\textcolor{red}{\textbf{[REVIEW]}}}{[REVIEW]}}
\newcommand{\thesis}{\texorpdfstring{\textcolor{blue}{\textbf{[THESIS]}}}{[THESIS]}}
\newcommand{\both}{\texorpdfstring{\textcolor{purple}{\textbf{[BOTH]}}}{[BOTH]}}

\begin{document}

\begin{titlepage}
\centering
\vspace*{3cm}

{\LARGE\bfseries Work Plan}\\[12pt]
{\Large Spectral Methods for Binary Neural Networks}

\vspace{2cm}

{\large Addressing Committee Review \& Research Milestones}

\vspace{1.5cm}

{\large February 2026}

\vspace{3cm}

\begin{tabular}{rl}
\textcolor{red}{\textbf{[REVIEW]}} & Task directly addresses committee feedback \\
\textcolor{blue}{\textbf{[THESIS]}} & Task from the proposed research program \\
\textcolor{purple}{\textbf{[BOTH]}} & Addresses review \emph{and} advances the research \\
\end{tabular}

\end{titlepage}

\tableofcontents
\newpage

% ══════════════════════════════════════════════════════════════════════════
\section{Executive Summary}

The committee's review identified five major and four moderate concerns. The single most critical action item is the \textbf{pilot spectral concentration experiment} --- the reviewer explicitly conditioned candidacy advancement on presenting this evidence. The second critical item is \textbf{at least one formal theoretical result}.

This work plan is organized into five phases, ordered so that each phase either (a)~addresses review concerns that gate subsequent work, or (b)~builds on infrastructure from the prior phase. The estimated total timeline is 10--12 months of focused work, with the pre-defense deliverables (Phases~1--2) completable in approximately 8--10 weeks.

\paragraph{Key principle.} We front-load the highest-risk, most informative experiments. If the pilot study shows no spectral concentration (Phase~1), we pivot early rather than investing months in an approach built on a false assumption.

% ══════════════════════════════════════════════════════════════════════════
\section{Phase 1: Pilot Study and Foundation (Weeks 1--4)}\label{sec:phase1}

\textit{Goal: Test the core assumption and build infrastructure. This phase produces the evidence the committee requires before the defense.}

\subsection{Task 1.1: Implement spectral analysis tooling \both}

\paragraph{What.} Build the \texttt{PilotAnalysis} module:
\begin{itemize}
    \item Fast Walsh--Hadamard transform for vectors of length $2^n$ (optimized NumPy/PyTorch implementation).
    \item Functions to compute: full Walsh--Fourier spectrum of a binary vector, spectral energy by degree $E(k) = \sum_{|S|=k} \hat{w}(S)^2$, effective spectral degree $k^*$ at the 95\% energy threshold, spectral entropy, and the per-row influence profile.
    \item Visualization: spectral energy heatmaps (layer $\times$ degree), histograms of $k^*$ across rows, comparison against random baselines.
\end{itemize}

\paragraph{Estimated effort.} 3--4 days. The FWHT is $\sim$20 lines of code; the analysis functions are straightforward.

\paragraph{Addresses.} Major Concern~1 (pilot experiment), Suggestion~1 (run it now).

\subsection{Task 1.2: Obtain and analyze binary weight matrices \both}

\paragraph{What.} Run the pilot study protocol described in the revised thesis (\S5):
\begin{enumerate}
    \item \textbf{Source A:} Download BitNet b1.58 2B4T open weights. Extract ternary weight matrices. Analyze as pseudo-Boolean functions.
    \item \textbf{Source B:} Load GPT-2 Small (124M). Binarize all weight matrices via $\text{sgn}(W)$. Analyze.
    \item \textbf{Source C:} (Deferred to Task~2.2 --- requires training a binary model first.)
    \item \textbf{Controls:} Random $\pm 1$ matrices of matched dimensions; random matrices with matched row-sum distribution; random linear threshold functions (LTFs) generated by $w(j) = \text{sgn}(\langle a, \text{bin}(j)\rangle)$ with $a \sim \mathcal{N}(0, I_n)$ where $n = \log_2 d$. The LTF control is critical: if trained weights look spectrally similar to random LTFs, the observed concentration reflects sign-activation structure rather than task-specific learning.
\end{enumerate}

For each matrix, compute the full spectral profile. Report the distribution of effective spectral degree $k^*$ across rows and layers.

\paragraph{Non-power-of-two dimensions.} Since BitNet b1.58 2B4T uses $d=2560$ and GPT-2 Small uses $d=768$, neither is a power of two. We handle this via (i) blockwise power-of-two submatrix analysis, and (ii) zero-padding to the next power of two, reporting both to quantify encoding artifacts.

\paragraph{Estimated effort.} 2--3 days for Sources~A and B plus controls. Estimated 1--6 CPU-hours depending on model size, threading, and permutation repeats.

\paragraph{Decision gate.} If $k^*$ is consistently close to $n$ (no concentration) across all sources: \textbf{stop and reassess.} Possible pivots:
\begin{itemize}
    \item Shift primary focus to Hadamard-structured layers (which do not depend on row-wise spectral sparsity).
    \item Investigate whether concentration emerges during spectral \emph{training} even if absent in conventionally trained weights.
    \item Reframe the thesis contribution around the theoretical framework and diagnostic tools rather than performance claims.
\end{itemize}

\paragraph{Addresses.} Major Concern~1 (the load-bearing assumption), Suggestion~1 (run the pilot now).


\subsection{Task 1.3: Column-indexing sensitivity analysis \review}

\paragraph{What.} Using the spectral data from Task~1.2, test whether the spectral profile is sensitive to column ordering:
\begin{enumerate}
    \item Verify (as theory predicts) that the degree-level energy profile $\sum_{|S|=k} \hat{w}(S)^2$ is invariant under bit permutations.
    \item Measure whether it changes significantly under 20 random column permutations.
    \item Compare standard binary encoding vs.\ Gray code encoding.
\end{enumerate}

\paragraph{Estimated effort.} 1 day (reuses infrastructure from Task~1.1).

\paragraph{Addresses.} Major Concern~2 (column-indexing arbitrariness), Suggestion~5.


\subsection{Task 1.4: Toy model --- Majority function \review}

\paragraph{What.} Implement and run the toy example from the revised thesis (\S4.8):
\begin{enumerate}
    \item Build a minimal \texttt{SpectralLinear} layer for $d = 8$ ($n = 3$).
    \item Train it to learn $\text{Maj}_3$ using spectral coefficients at degree $k = 1, 2, 3$.
    \item Verify convergence to the known Fourier spectrum.
    \item Extend to $\text{Maj}_7$ ($d = 128$, $n = 7$) and the tribes function.
    \item Log all spectral diagnostics; visualize training dynamics.
\end{enumerate}

\paragraph{Estimated effort.} 2--3 days.

\paragraph{Addresses.} Suggestion~6 (toy model), and serves as a sanity check for the \texttt{SpectralLinear} implementation before scaling.

\subsection{Phase 1 Deliverables}

\begin{enumerate}
    \item Spectral concentration results for BitNet and binarized GPT-2 weight matrices (figures + summary statistics).
    \item Column-indexing sensitivity report.
    \item Working toy-model demonstration with training curves and spectral diagnostics.
    \item Go/no-go decision on the row-wise spectral parameterization.
\end{enumerate}


% ══════════════════════════════════════════════════════════════════════════
\section{Phase 2: Theory and Core Implementation (Weeks 5--10)}\label{sec:phase2}

\textit{Goal: Produce at least one formal theoretical result and build the core \texttt{SpectralLinear} and \texttt{HadamardLinear} modules. Combined with Phase~1, this constitutes the defense-ready package.}

\subsection{Task 2.1: Strengthen theoretical results \review}

The revised thesis already includes two preliminary results (the degree-$k$ approximation bound and the normalization propagation proposition). The committee asks for ``at least one formal theoretical result.'' We aim for three contributions, in decreasing order of confidence:

\paragraph{Result A (high confidence): Tighten the approximation bound.}
The current Theorem~3.1 in the revised thesis bounds the error of degree-$k$ truncation in terms of total influence. Strengthen this by:
\begin{itemize}
    \item Proving a tighter bound for the specific case where $f$ is a linear threshold function, using quantitative noise sensitivity bounds for LTFs (e.g., from Harsha et al.).
    \item Deriving an explicit bound on the number of bit flips introduced by the sign projection as a function of $k$ and the spectral gap.
\end{itemize}

\paragraph{Result B (medium confidence): Spectral-domain gradient characterization.}
Prove that the gradient of the cross-entropy loss with respect to spectral coefficients $\hat{w}(S)$ decomposes as:
\[
\frac{\partial \mathcal{L}}{\partial \hat{w}_i(S)} = \sum_j \frac{\partial \mathcal{L}}{\partial w_i(j)} \cdot \chi_S(\text{bin}(j)),
\]
which is itself a Walsh--Fourier coefficient of the spatial gradient. Note that this identity alone is just a change of basis and does not imply any advantage. The formal goal is to show that \emph{combined with degree truncation or $L_1$/degree-penalty regularization in the spectral domain}, gradient descent converges faster to low-complexity Boolean functions than weight-domain STE. This is a formal version of the ``better inductive bias'' claim; without the sparsity-inducing component, no separation should be expected.

\paragraph{Result C (aspirational): Separation result.}
Show that there exists a family of Boolean functions for which spectral parameterization at degree $k = O(\log n)$ converges to the optimum in polynomial time, while weight-domain STE requires exponential time. The candidate family: random $k$-juntas (functions depending on only $k$ of $n$ variables), for which the Kushilevitz--Mansour algorithm gives polynomial-time learning but brute-force search over $\pm 1$ weights is exponential.

\paragraph{Estimated effort.} 2--3 weeks. Result A is a relatively straightforward extension of known bounds. Result B requires careful but standard calculus of the WHT. Result C is genuinely hard and may remain a conjecture.

\paragraph{Pivot plan for Result B.} If a clean formal proof of the gradient advantage cannot be obtained within the allocated time, we pivot to a rigorous empirical result: measure gradient variance, spectral sparsity, and convergence speed in both spectral and weight domains across multiple random seeds, problem scales, and regularization strengths. Present the comparison with confidence intervals and statistical tests. This produces a ``rigorous empirical result'' that satisfies the committee's requirement for formal analysis even if a closed-form proof eludes us.

\paragraph{Addresses.} Major Concern~5 (missing theory).


\subsection{Task 2.2: Implement \texttt{SpectralLinear} \thesis}

\paragraph{What.} Production-quality PyTorch module:
\begin{itemize}
    \item Parameters: spectral coefficient tensor of shape $[d_\text{out}, \binom{n}{\leq k}]$.
    \item Precomputed: Walsh basis matrix of shape $[\binom{n}{\leq k}, d_\text{in}]$ (rows of $H_n$ corresponding to monomials of degree $\leq k$).
    \item Forward pass: $W_\text{cont} = \hat{w} \cdot \Phi^\top$; $W_\text{bin} = \tanh(\beta \cdot W_\text{cont})$ (or $\text{sgn}$ with STE); output $= W_\text{bin} \cdot x$.
    \item Backward pass: autograd handles everything (STE through sign; exact through linear WHT map).
    \item Spectral norm projection as a post-step hook.
\end{itemize}

\paragraph{Estimated effort.} 1 week.

\paragraph{Validation.} Verify on the toy tasks from Task~1.4 that the module produces identical results to the prototype. Benchmark forward/backward speed against \texttt{nn.Linear}.


\subsection{Task 2.3: Implement \texttt{HadamardLinear} \thesis}

\paragraph{What.} PyTorch module implementing the butterfly factorization:
\begin{itemize}
    \item Parameters: $L_H$ sign vectors $s_\ell \in \{-1, +1\}^d$; $L_H$ soft permutation matrices (Sinkhorn-parameterized during training, hard at inference).
    \item Forward: cascade of FWHT $\to$ sign-flip $\to$ permute, with explicit $d^{-L_H/2}$ normalization.
    \item FWHT implementation using the standard $O(d \log d)$ butterfly algorithm.
\end{itemize}

\paragraph{Estimated effort.} 1 week.

\paragraph{Addresses.} Moderate Concern~7 (Sinkhorn leakage) --- implement hard-rounding ablation from the start.


\subsection{Task 2.4: Train a baseline binary Transformer (STE) \thesis}

\paragraph{What.} Train Baseline~2 (standard binary STE, no spectral parameterization) on WikiText-103 at $d = 256$, $L = 4$. This serves three purposes:
\begin{enumerate}
    \item Provides the binary weight matrices needed for Source~C of the pilot study (completing Task~1.2).
    \item Establishes the baseline perplexity to beat.
    \item Validates the training infrastructure (data loading, evaluation, logging) before adding spectral components.
\end{enumerate}

\paragraph{Estimated effort.} 3--5 days (training time dominates; implementation reuses standard Transformer code with sign quantization).


\subsection{Phase 2 Deliverables}

\begin{enumerate}
    \item At least one (ideally two) formal theoretical results with complete proofs.
    \item Working, tested \texttt{SpectralLinear} and \texttt{HadamardLinear} modules.
    \item STE baseline results on WikiText-103 ($d = 256$, $L = 4$).
    \item Completed pilot study including Source~C (STE-trained binary weights).
    \item \textbf{Defense-ready package:} revised thesis + pilot study results + theoretical results.
\end{enumerate}


% ══════════════════════════════════════════════════════════════════════════
\section{Phase 3: Core Experiments (Weeks 11--20)}\label{sec:phase3}

\textit{Goal: Train and evaluate the Spectral Binary Transformer. Run the primary ablation studies.}

\subsection{Task 3.1: Implement the full SBT architecture \thesis}

\paragraph{What.} Assemble \texttt{SpectralTransformerBlock} from the modules built in Phase~2:
\begin{itemize}
    \item \texttt{SpectralLinear} for $Q$, $K$, $V$, output projection, and FFN layers.
    \item Integer attention with fixed $1/\sqrt{d_k}$ scaling (Option~A --- the recommended starting point).
    \item RMSNorm or LayerNorm (at appropriate precision).
    \item Residual connections.
    \item Spectral diagnostics logging (effective degree, entropy, rounding error per layer per $N$ steps).
\end{itemize}

\paragraph{Estimated effort.} 1 week.


\subsection{Task 3.2: Small-scale SBT training runs \thesis}

\paragraph{What.} Train the SBT at the smallest configuration ($d = 256$, $L = 4$, $k = 3$) on WikiText-103:
\begin{enumerate}
    \item Establish that training converges (loss decreases; spectral diagnostics are sensible).
    \item Compare against Baseline~2 (STE) and Baseline~4 (random frozen binary).
    \item Sweep learning rate, temperature schedule $\beta_t$, and regularization strength $\lambda$.
\end{enumerate}

\paragraph{Estimated effort.} 1--2 weeks (many short training runs).

\paragraph{Decision gate.} If SBT at $d = 256$ cannot outperform random binary weights after hyperparameter tuning: investigate whether the issue is (a)~the spectral degree $k$, (b)~the projection method, or (c)~something more fundamental. Run ablations A1 and A2 at this scale before scaling up.


\subsection{Task 3.3: Primary ablation studies \thesis}

Run ablations A1--A6 at the small scale ($d = 256$, $L = 4$):

\begin{center}
\begin{tabular}{lll}
\toprule
\textbf{Ablation} & \textbf{Variable} & \textbf{Settings} \\
\midrule
A1 & Spectral degree $k$ & $1, 2, 3, 4, \ldots, 8$ \\
A2 & Projection method & Hard sign, randomized, spectral norm \\
A3 & Structured vs.\ unstructured & SpectralLinear vs.\ HadamardLinear \\
A4 & Spectral regularization & $\lambda \in \{0, 10^{-4}, 10^{-3}, 10^{-2}\}$; $\mu \in \{0, 10^{-3}\}$ \\
A5 & Activation precision & Binary / 8-bit / full precision \\
A6 & Embedding strategy & Learned FP / learned binary / Hadamard \\
\bottomrule
\end{tabular}
\end{center}

\paragraph{Estimated effort.} 2--3 weeks (many runs, but each is fast at this scale).

\paragraph{Key outputs.}
\begin{itemize}
    \item Optimal $k^*$ for the spectral degree (A1). This is the most important result.
    \item Whether projection method matters significantly (A2).
    \item Whether Hadamard structure helps or hurts (A3).
    \item Whether spectral regularization improves convergence (A4).
\end{itemize}


\subsection{Task 3.4: Scale-up experiments \thesis}

Using the best configuration from Task~3.3, train at larger scales:
\begin{enumerate}
    \item $d = 512$, $L = 8$ on WikiText-103.
    \item $d = 1024$, $L = 12$ on WikiText-103 (and/or Pile subset, if compute permits).
    \item Compare against all four baselines at each scale.
\end{enumerate}

\paragraph{Estimated effort.} 2--3 weeks (longer training runs; fewer configurations).


\subsection{Task 3.5: Attention mechanism comparison \review}

\paragraph{What.} At the best-performing scale from Task~3.4, compare the three attention variants:
\begin{enumerate}
    \item Option A: Integer attention with fixed scaling (already implemented).
    \item Option B: Hadamard attention (implement from the equations in the revised thesis \S4.5).
    \item Option C: Linear attention with binary feature maps (implement binary LSH kernel).
\end{enumerate}

Measure perplexity, training stability, and inference arithmetic profile (what fraction of operations are binary).

\paragraph{Estimated effort.} 1--2 weeks.

\paragraph{Addresses.} Moderate Concern~8 (attention mechanisms underdeveloped).


\subsection{Phase 3 Deliverables}

\begin{enumerate}
    \item Complete ablation results at small scale (figures, tables).
    \item Perplexity results at $d \in \{256, 512, 1024\}$ against all baselines.
    \item Attention mechanism comparison.
    \item Spectral diagnostic plots over training for all configurations.
    \item Determination of whether ``strong,'' ``moderate,'' or ``diagnostic'' success criterion is met.
\end{enumerate}


% ══════════════════════════════════════════════════════════════════════════
\section{Phase 4: Extended Experiments and Analysis (Weeks 21--30)}\label{sec:phase4}

\textit{Goal: Run the remaining ablations, perform the zero-shot evaluation, and develop the inference pipeline analysis.}

\subsection{Task 4.1: Additional ablations (A7, A8) \review}

\begin{itemize}
    \item \textbf{A7: Sinkhorn rounding sensitivity.} For Hadamard-structured layers, measure performance when soft Sinkhorn permutations are rounded to hard permutations at various points during training (early, mid, late, only at end).
    \item \textbf{A8: Column encoding.} Standard binary vs.\ Gray code vs.\ learned permutation.
\end{itemize}

\paragraph{Estimated effort.} 1 week.

\paragraph{Addresses.} Moderate Concerns~2 (column indexing) and 7 (Sinkhorn leakage).


\subsection{Task 4.2: Zero-shot evaluation \thesis}

\paragraph{What.} Evaluate the best SBT model on downstream benchmarks via lm-evaluation-harness:
\begin{itemize}
    \item LAMBADA (cloze completion)
    \item HellaSwag (commonsense reasoning)
    \item ARC-Easy and ARC-Challenge (science QA)
    \item PIQA (physical intuition)
    \item WinoGrande (coreference)
\end{itemize}

Compare against all baselines. This is the standard evaluation expected by the binary network community.

\paragraph{Estimated effort.} 3--5 days (evaluation is fast; the work is in setting up lm-eval-harness integration).


\subsection{Task 4.3: Inference pipeline analysis \review}

\paragraph{What.} Produce a complete accounting of the arithmetic operations in the SBT inference pipeline:
\begin{itemize}
    \item For each layer: count binary operations (additions/subtractions via XNOR+popcount), integer operations, and any remaining floating-point operations.
    \item Identify exactly where floating-point leaks in: softmax, normalization, embeddings.
    \item Quantify: what fraction of total operations are strictly binary? How does this compare to BitNet?
    \item Prototype a pure-integer inference path (if possible) using fixed-point arithmetic for the non-binary components.
\end{itemize}

\paragraph{Estimated effort.} 1--2 weeks.

\paragraph{Addresses.} Major Concern~3 (scope of ``purely binary'' claim), Moderate Concern~6 (scaling arithmetic), Suggestion~3 (tighten claims or qualify).


\subsection{Task 4.4: Multi-layer composition analysis \review}

\paragraph{What.} Empirically study how spectral structure evolves through the network:
\begin{itemize}
    \item At each layer, compute the effective spectral degree of the \emph{activations} (after binarization), not just the weights.
    \item Track noise stability $\text{Stab}_\rho$ across layers. Does it decay with depth, and under what conditions? The composition identity in \S2.8 of the thesis provides a framework for interpreting the results.
    \item Monitor the mutual information diagnostics from the thesis (\S4.1): the empirical output entropy $H(\text{sgn}(W_\ell y_{\ell-1}))$ per layer (estimated over minibatches) and the per-row total influence $\mathbf{I}[f_i]$. If $H(\text{sgn}(Wx))$ drops below $\log_2 V$ at any layer, the sign bottleneck is binding.
    \item Correlate layer-wise spectral properties with gradient magnitudes.
\end{itemize}

\paragraph{Estimated effort.} 1 week.

\paragraph{Addresses.} Major Concern~4 (activation precision hand-waved), the new Open Problem~3 (multi-layer composition).


\subsection{Phase 4 Deliverables}

\begin{enumerate}
    \item Zero-shot benchmark results table.
    \item Complete inference arithmetic breakdown (table + discussion).
    \item Layer-wise spectral composition analysis (figures).
    \item All remaining ablation results (A7, A8).
\end{enumerate}


% ══════════════════════════════════════════════════════════════════════════
\section{Phase 5: Writing and Defense (Weeks 31--40)}\label{sec:phase5}

\subsection{Task 5.1: Thesis writing}

Expand the proposal into a full thesis. The revised proposal already contains substantial text for Chapters~1--4. The remaining writing:

\begin{itemize}
    \item \textbf{Chapter 5: Pilot Study Results.} Present the spectral concentration findings from Phase~1 with figures and analysis.
    \item \textbf{Chapter 6: Experimental Results.} All results from Phases~3--4, organized by ablation and scale.
    \item \textbf{Chapter 7: Analysis and Discussion.} Inference pipeline analysis, multi-layer composition study, and honest assessment of which success criterion was met.
    \item \textbf{Chapter 8: Conclusion and Future Work.} What was learned, what remains open, and concrete next steps.
\end{itemize}

\paragraph{Estimated effort.} 4--6 weeks.


\subsection{Task 5.2: Open-source release}

Package and release:
\begin{itemize}
    \item \texttt{spectral-binary-nn}: Python package with \texttt{SpectralLinear}, \texttt{HadamardLinear}, \texttt{SpectralDiagnostics}, and \texttt{PilotAnalysis}.
    \item Pretrained SBT checkpoints at each scale.
    \item Reproducibility scripts for all experiments.
\end{itemize}

\paragraph{Estimated effort.} 1 week (code cleanup, documentation, tests).


% ══════════════════════════════════════════════════════════════════════════
\section{Summary: Review Concern Traceability}\label{sec:traceability}

The following table maps each committee concern to where it is addressed in both the revised thesis and this work plan.

\begin{center}
\small
\begin{longtable}{p{0.28\textwidth} p{0.30\textwidth} p{0.32\textwidth}}
\toprule
\textbf{Committee Concern} & \textbf{Revised Thesis} & \textbf{Work Plan Task} \\
\midrule
\endfirsthead
\toprule
\textbf{Committee Concern} & \textbf{Revised Thesis} & \textbf{Work Plan Task} \\
\midrule
\endhead
\textbf{Major 1:} Core analogy may not hold & \S2.4 expanded remark; new \S5 (Pilot Study) & Tasks 1.1, 1.2 (pilot study) \\
\midrule
\textbf{Major 2:} Column-indexing arbitrary & New \S4.2.2; invariance argument & Tasks 1.3, 4.1 (ablation A8) \\
\midrule
\textbf{Major 3:} ``Purely binary'' oversold & New ``Scope and claims'' paragraph in \S1.3; precise three-part claim & Task 4.3 (inference analysis) \\
\midrule
\textbf{Major 4:} Activation precision hand-waved & New \S2.8 (multi-layer composition); composition identity (no longer a numbered proposition) & Task 4.4 (composition analysis); Ablation A5 \\
\midrule
\textbf{Major 5:} Missing theory & New \S3 (Theoretical Contributions): Prop.~3.1, Prop.~3.2, Conjectures 3.5--3.6 & Task 2.1 (strengthen theory) \\
\midrule
\textbf{Moderate 6:} Scaling arithmetic & Prop.~3.2 \& Remark 3.3 (normalization propagation) & Task 4.3 (inference pipeline) \\
\midrule
\textbf{Moderate 7:} Sinkhorn leakage & Acknowledged in \S4.3; new Ablation A7 & Tasks 2.3, 4.1 (ablation A7) \\
\midrule
\textbf{Moderate 8:} Attention underdeveloped & \S4.4 expanded with equations, ranking, and rationale & Task 3.5 (comparison) \\
\midrule
\textbf{Moderate 9:} Bibliography issues & WHNO cite key fixed (\texttt{cavallazzi2025whno}); Merity date corrected; Alman arXiv corrected; foundational BNN references added & (Done in revision) \\
\midrule
\textbf{Suggestion:} Pilot experiment & New \S5 with full protocol & Tasks 1.1--1.2 \\
\midrule
\textbf{Suggestion:} Theory section & New \S3 & Task 2.1 \\
\midrule
\textbf{Suggestion:} Tighten binary claims & \S1.3 ``Scope and claims'' & Task 4.3 \\
\midrule
\textbf{Suggestion:} Commit to attention & \S4.4 with ranked options & Task 3.5 \\
\midrule
\textbf{Suggestion:} Column-indexing & \S4.2.2 & Tasks 1.3, 4.1 \\
\midrule
\textbf{Suggestion:} Toy model & New \S4.7 (Majority function) & Task 1.4 \\
\bottomrule
\end{longtable}
\end{center}


% ══════════════════════════════════════════════════════════════════════════
\section{Timeline Overview}\label{sec:timeline}

\begin{center}
\begin{tabular}{llll}
\toprule
\textbf{Phase} & \textbf{Weeks} & \textbf{Key Milestone} & \textbf{Gate?} \\
\midrule
\rowcolor{phase1}
1: Pilot \& Foundation & 1--4 & Spectral concentration results & Yes \\
\rowcolor{phase2}
2: Theory \& Core Impl. & 5--10 & Defense-ready package & --- \\
\midrule
\multicolumn{4}{c}{\textit{--- Proposal Defense ---}} \\
\midrule
\rowcolor{phase3}
3: Core Experiments & 11--20 & SBT results vs.\ baselines & Yes \\
\rowcolor{phase4}
4: Extended \& Analysis & 21--30 & Zero-shot eval; inference analysis & --- \\
\rowcolor{phase5}
5: Writing \& Defense & 31--40 & Thesis submission & --- \\
\bottomrule
\end{tabular}
\end{center}

\paragraph{Decision gates.}
\begin{itemize}
    \item \textbf{End of Phase~1:} If no spectral concentration observed, pivot to Hadamard-structured approach and/or reframe contributions. Do not proceed to Phase~3 with the row-wise spectral parameterization.
    \item \textbf{End of Phase~3 (Task~3.2):} If SBT at small scale cannot outperform random binary weights, diagnose and fix before scaling up.
\end{itemize}


% ══════════════════════════════════════════════════════════════════════════
\section{Resource Requirements}\label{sec:resources}

\begin{center}
\begin{tabular}{lll}
\toprule
\textbf{Phase} & \textbf{Compute} & \textbf{Notes} \\
\midrule
1 & CPU only & Spectral analysis; toy models \\
2 & 1 GPU (A100 or equiv.) & STE baseline training \\
3 & 2--4 GPUs & Multiple training runs at $d \leq 1024$ \\
4 & 1--2 GPUs & Evaluation; analysis \\
5 & Minimal & Writing \\
\bottomrule
\end{tabular}
\end{center}

Total estimated GPU-hours: 2,000--5,000 (modest by current standards; all models are small by design).


% ══════════════════════════════════════════════════════════════════════════
\section{Immediate Next Actions (This Week)}\label{sec:immediate}

\begin{enumerate}
    \item \textbf{Monday--Tuesday:} Implement the FWHT and spectral analysis functions (Task~1.1).
    \item \textbf{Wednesday:} Download BitNet b1.58 2B4T weights; binarize GPT-2 Small weights. Run full spectral analysis on both (Task~1.2).
    \item \textbf{Thursday:} Column-indexing sensitivity tests (Task~1.3). Generate all pilot study figures.
    \item \textbf{Friday:} Implement and run the Majority function toy model (Task~1.4). Review Phase~1 results and make the go/no-go decision.
\end{enumerate}

If Phase~1 results are positive, begin Task~2.1 (theory) and Task~2.2 (SpectralLinear implementation) the following Monday.

\end{document}
